(1) 球面积 \(A=4\pi R^2\)，辐射净流出为 \(\epsilon\sigma A(T^4-T_a^4)\)，线性对流流出为 \(hA(T-T_a)\)。以均匀球温为唯一状态变量，集总能量方程是 \(C\dot T=P-\epsilon\sigma A(T^4-T_a^4)-hA(T-T_a)\)，仅在球内热均衡比外部散热快时成立。
(2) 定义 \(F(T)=\epsilon\sigma A(T^4-T_a^4)+hA(T-T_a)\)。对 \(T>0\)，\(F'(T)=4\epsilon\sigma AT^3+hA>0\)（至少一项正），且 \(F(T_a)=0\)、\(F(T)\to\infty\)，所以每个 \(P>0\) 对应唯一 \(T_*>T_a\)。当 \(T>T_*\)，\(\dot T<0\)，当 \(T<T_*\)，\(\dot T>0\)，故平衡渐近稳定。
(3) 对 \(T=T_*+\delta T\) 线性展开，\(C\dot{\delta T}=-[4\epsilon\sigma AT_*^3+hA]\delta T+O(\delta T^2)\)，因此 \(\delta T(t)=\delta T(0)e^{-t/\tau}\)，\(\tau=C/[A(4\epsilon\sigma T_*^3+h)]\)。该时间尺度用总热容除以平衡点处的微分热导而不是除以总辐射功率。
(4) 对小 \(P\)，\(T_*=T_a+\Delta T\)，一阶为 \(P\simeq A(4\epsilon\sigma T_a^3+h)\Delta T\)，故 \(\Delta T\simeq P/[A(4\epsilon\sigma T_a^3+h)]\)。若只向零温背景辐射且无输入，则 \(C\dot T=-\epsilon\sigma AT^4\)，积分得 \(T(t)=[T(0)^{-3}+3\epsilon\sigma At/C]^{-1/3}\)，呈代数而非指数衰减。
